\documentclass[../DoAn.tex]{subfiles}
\begin{document}

\begin{center}
    \Large{\textbf{ABSTRACT}}\\
\end{center}
\vspace{1cm}

Modern mobile devices can acquire GNSS signals, inertial measurements and camera frames at high frequency. During road tracking, however, GNSS may become weak or unavailable because of occlusion, dense urban environments or device limitations. This thesis develops an Android application that explores how a phone camera and inertial sensors can help update the user location directly on a map when GNSS quality degrades. From this objective, the system is organized into three problems: GNSS Viewer, Optical Flow View and LiveRouting that combines GNSS with Optical Flow/IMU. In GNSS Viewer, the application collects location updates, satellite status and raw GNSS measurements, then resolves satellite positions using Real GNSS PVT, IGS Broadcast Ephemeris, CelesTrak TLE/SGP4 and an azimuth/elevation approximation fallback. The results are displayed on a two-dimensional map, an OpenGL 3D Earth model with day/night rendering, Moon, Sun, Orekit TAI-UTC time data and an AR view that overlays satellites on the camera stream. In Optical Flow View, the application implements KLT and Farneback with OpenCV for real-time camera analysis, supports ROI selection, object tracking, video recording, stored analysis sessions and metrics such as FPS, processing time, active vectors and confidence. For uploaded videos, the system supports both the custom FastAPI server and an external team backend; the ONNX RAFT pipeline processes videos as jobs, supports chunked upload/download, cancellation, result retrieval and Cutie-based object ROI. In LiveRouting, the application tracks the route on the map; when GNSS is not reliable, camera assist is activated so Optical Flow and IMU can provide a pragmatic visual odometry estimate, confidence-based route map matching and separate path segments for testing and later inspection.

The current LiveRouting design is limited to vehicle movement. When GNSS becomes unreliable, Optical Flow metrics are converted into a pragmatic visual odometry estimate containing speed, heading change and confidence; IMU yaw rate stabilizes the heading; and the current route is used for confidence-based map matching instead of hard snapping. The estimator keeps an internal position uncertainty that grows during GNSS outage, blends the predicted pose back to the route only when distance, heading and continuity are reliable, and triggers route refresh when the estimated trajectory deviates from the active route.

\end{document}
